\begin{problem}[33届物理竞赛复赛]{63}
    质子是由更小的所谓“部分子”构成的。欧洲大型强子对撞机（LHC）是高能质子-质子对撞机，质子束内单个质子能量为 $E = 7.0 \, \text{TeV}$ （ $1\, \text{TeV} = 10^3\, \text{GeV} = 10^{12}\, \text{eV}$ ），两束能量相同的质子相向而行对撞碎裂，其中相撞的两个部分子 a、b 相互作用湮灭产生一个新粒子。设部分子 a、b 的动能在质子能量中所占的比值分别为 $x_a$ 、 $x_b$ ，且远大于其静能。

    \begin{subproblem}
        假设两个部分子 a、b 对撞湮灭产生了一个静质量为 $m_{s}=1.0\;TeV/c^{2}$ 的新粒子 S，求 $x_{a}$ 和 $x_{b}$ 的乘积 $x_{a}x_{b}$ ；
    \end{subproblem}
    \begin{subproblem}
        假设新粒子 S 产生后衰变到两个光子，在新粒子 S 静止的参考系中，求两光子的频率；
    \end{subproblem}
    \begin{subproblem}
        假设新粒子 S 产生后在其静止坐标系中衰变到两个质量为 $m_{A} = 1.0 \, GeV/c^{2}$ 的轻粒子 A，每个轻粒子 A 再衰变到两个同频率的光子，求在这个坐标系中这两个光子动量之间的夹角。
    \end{subproblem}

    已知： $\sin\alpha\approx\alpha$ ，当 $\alpha<<1$ ；普朗克常量 $h=6.63\times10^{-34}J\cdot s$ ，电子电荷量绝对值 $e=1.60\times10^{-19}C$ 。
\end{problem}